(1) 将微分方程直接积分 $A_2(L)=i\kappa A_1^2\int_0^Le^{i\Delta kz}dz$，故 $\boxed{A_2(L)=i\kappa A_1^2Le^{i\Delta kL/2}\operatorname{sinc}(\Delta kL/2)}$，其中 $\operatorname{sinc}x=\sin x/x$，$x=0$ 处连续取1。

(2) 输出强度与 $|A_2|^2$ 成正比，所以 $\boxed{I_2(L)\propto|\kappa|^2|A_1|^4L^2\operatorname{sinc}^2(\Delta kL/2)}$。相位失配决定沿程微小源项是否同相累积；整体因子 $e^{i\Delta kL/2}$ 只改变二次谐波相位，不改变输出强度。

(3) 当 $\Delta k=0$，$|A_2|\propto L$、$I_2\propto L^2$。若 $\Delta k\ne0$，第一处输出相消为 $|\Delta k|L/2=\pi$，即 $\boxed{L=2\pi/|\Delta k|}$；此时不同晶体位置辐射的相位源项在振幅上互相抵消。

(4) 弱耗尽要求转换到二次谐波的能量远小于入射基频能量，因此将 $A_1$ 视为固定。高效率时必须联立两频场，Manley–Rowe关系保证每产生一个 $2\omega$ 光子消耗两个 $\omega$ 光子。周期性翻转 $\chi^{(2)}$ 可补偿积累的相位失配，实现准相位匹配而不是凭空创造能量。
